%!TEX root = ./../../springer-mono-format/book.tex
%!TEX root = ./../../thesis-trier.tex

\IN{completely distinguishable!automaton}
An automaton $\semiautomat{A}{Q}{\Sigma}{\delta}$ is called \emph{completely subset-pair distinguishable} (or just \emph{completely distinguishable} for short) 
if every two distinct non-empty and non-singleton subsets $A,B \subseteq Q$ are distinguishable in~$\mathcal A$.

Recall that $L(\mathcal P(\mathcal A)_{Q, F}) = \Syn(\mathcal A)$
with $F = \{\{q\} \mid q \in Q\}$, see~\cref{sync_max:sec:power_automat_for_SynA}. 
By \crefitem{prelim:thm:char_minaut}{thm:minaut:item:statesdist},
distinguishability of pairs of distinct states is a necessary condition 
for minimality.
So, applying this to $\mathcal P(\mathcal A)_{Q, F}$,
complete distinguishability
is a necessary condition for the minimal automaton
of $\Syn(\mathcal A)$ to have $2^{|Q|} - |Q|$ states.

First, in \cref{sync_max:thm:cerny_compl_dist}, as an example,
we prove that the \Cerny automata from \cref{intro:fig:cerny}
are completely distinguishable (this also follows from \cref{sync_max:thm:circular_sync_max} given later,
but as an example we give a self-contained proof here).
Then, in \cref{sync_max:prop:dist_sync},
we show that completely distinguishable automata
with at least three states are synchronizing.
\Cref{sync_max:thm:char_distinguishable} collects various equivalent characterizations,
and \cref{sync_max:cor:sufficient-compl-dist,cor:compl-dist-transitive}
are two consequences of these characterizations.
\Cref{sync_max:ex:compl-dist-non-cr,sync_max:ex:compl-dist-non-cr2,syncmax:ex:compl-dist-not-SCC} 
give completely distinguishable automata that are not completely reachable.
\Cref{syncmax:prop:binary-rank-two-letter} 
proves that in binary completely distinguishable automata where one letter has rank two,
the other letter cyclically permutes the states.
\cref{sync_max:prop:dist-n-1-set} lists an equivalence for a set missing one state being
distinguishable from the whole state set (a condition that appears in \cref{sync_max:thm:char_distinguishable}
for each set missing one state).

\begin{theorem}
	\label{sync_max:thm:cerny_compl_dist}
	The \Cerny automata $\mathcal C_n = ([n], \{a,b\}, \delta_n)$ from \cref{intro:fig:cerny}
	are completely distinguishable.
\end{theorem}
\begin{proof}
For $i \in [n]$ 
set $u_i = b^{n - i} a b^i$.
Then, for $j \in [n]$,
\begin{equation} \label{sync_max:thm:cerny_compl_dist:eqn}
\applytrans{\delta_n}{j}{u_i} = \begin{cases}
	(j + 1 ) \bmod n & \text{if } i = j, \\
	j             & \text{if } i \ne j.	
\end{cases}
\end{equation}
Now, let $A, B \subseteq [n]$ be distinct with $|A|, |B| \ge 2$. %two distinct non-empty and non-singleton subsets. 
Without loss of generality, 
suppose $i \in A \setminus B$. Set $u = u_{i+1} \cdots u_0 u_1 \cdots u_{i-1}$.
Then $\delta(i, u) = i$, $\applytrans{\delta}{B}{u} \subseteq \applytrans{\delta}{[n] \setminus \{i\}}{u} = \{ i - 1 \}$
and $\applytrans{\delta}{A}{u} = \{i,i-1\}$. So,~$u$
distinguishes~$A$ from~$B$.
\end{proof} % then u = ... disnigushes a from B as ...

\crefitem{sync_max:cor:two-sets-containing-2set}{sync_max:cor:two-sets-containing-2set:item:kl} yields
that completely distinguishable automata with at least three states are synchronizing.

\begin{proposition}
	\label{sync_max:prop:dist_sync}
	Let $\semiautomat{A}{Q}{\Sigma}{\delta}$ 
	be a completely distinguishable automaton
	with $|Q| \ge 3$. Then $\mathcal A$
	is synchronizing.
\end{proposition} 
\begin{proof}
	This follows from \crefitem{sync_max:cor:two-sets-containing-2set}{sync_max:cor:two-sets-containing-2set:item:kl},
	since every $2$-set is distinguishable from every $3$-set.
\end{proof}

Since the empty set is distinguishable from every other set if and only if 
the automaton is synchronizing, by \cref{sync_max:prop:dist_sync}
an automaton with at least three states is completely distinguishable if every pair 
of non-singleton subsets is distinguishable.

\begin{remark}\label{syncmax:rem:two-state-compl-dist}
	An automaton with at most two states admits no pair of distinct non-empty
	and non-singleton subsets of states. Thus, it is
	completely distinguishable, even if it is non-synchronizing.	
\end{remark}

Next, we formulate various equivalent characterizations.

\begin{theorem}
	\label{sync_max:thm:char_distinguishable}
	Let $\semiautomat{A}{Q}{\Sigma}{\delta}$ be an $n$-state automaton
	with $n \ge 3$.
	The following are equivalent:
	\begin{enumerate}
		\item \label{dist:item:all_pairs}
		$\mathcal A$ is completely distinguishable.
		
		\item \label{dist:item:subsets} 
		For every $A, B \subseteq Q$
		with $A \subsetneq B$ and $A$ non-empty
		there exists $u \in \Sigma^*$
		such that $1 = |\applytrans{\delta}{A}{u}| < |\applytrans{\delta}{B}{u}|$.
		
		\item \label{dist:item:subset_single_point} For every non-empty $A \subseteq Q$ and $q \in Q \setminus A$ there exists $u \in \Sigma^*$ 
		such that $\applytrans{\delta}{q}{u} \notin \applytrans{\delta}{A}{u}$
		and $|\applytrans{\delta}{A}{u}| = 1$.		
						
		\item \label{dist:item:single_point} For every $q \in Q$ there exists $u \in \Sigma^*$ such that
		$1 = |\applytrans{\delta}{Q \setminus \{q\}}{u}|$ and $\applytrans{\delta}{q}{u} \notin \applytrans{\delta}{Q \setminus \{q\}}{u}$, i.e., $Q$ is distinguishable from every $(n-1)$-set.
		
		\item \label{dist:item:one-two-set} For two distinct non-empty and non-singleton subsets
		there exists a word mapping one to a singleton and the other to a $2$-set.
					
		\item \label{dist:item:pqr-r-one-lane} For distinct $p,q,r \in Q$ there exists $u \in \Sigma^*$ 
		such that $\delta(p, u) = \delta(q, u)$ and 
		\[
		\delta^{-1}(\{\delta(r,u)\}, u) = \{r\},
		\] i.e.,
		$\delta(s, u) \ne \delta(r, u)$ for all $s \in Q \setminus \{r\}$.
		
		\item \label{dist:item:compress-two-sets} For every $A, B \subseteq Q$ with $\min\{|A|, |B|\} \ge 2$
		and $A \subsetneq B$ there exists $u \in \Sigma^*$ such that $|\delta(A, u)| < |\delta(B, u)| < |B|$.		
				
		\item \label{dist:item:compress} For every $A \subseteq Q$
		with $|A| \ge 2$ and $p \in Q \setminus A$
		there exists $u \in \Sigma^*$
		such that $\applytrans{\delta}{p}{u} \notin \applytrans{\delta}{A}{u}$
		and $|\applytrans{\delta}{A}{u}| < |A|$. 
				
		\item \label{dist:item:dcps-nondisjoint} For every distinct non-empty and non-singleton $A, B \subseteq Q$ with $A \cap B \ne \emptyset$
		there exists $u \in \Sigma^*$
		with $3 \le |\applytrans{\delta}{A}{u}| + |\applytrans{\delta}{B}{u}| < |A| + |B|$
		and $\applytrans{\delta}{A}{u} \ne \applytrans{\delta}{B}{u}$.		
		
		\item \label{dist:item:dcps} For every distinct non-empty and non-singleton $A, B \subseteq Q$  
		there exists $u \in \Sigma^*$
		with $3 \le |\applytrans{\delta}{A}{u}| + |\applytrans{\delta}{B}{u}| < |A| + |B|$
		and $\applytrans{\delta}{A}{u} \ne \applytrans{\delta}{B}{u}$.		
	\end{enumerate}
\end{theorem}
\begin{proof}
	$\ref{dist:item:all_pairs} \Rightarrow \ref{dist:item:subsets}$:
	By assumption there exists $u \in \Sigma^*$
	distinguishing $A$ from $B$.
	As $A \subsetneq B$
	we must have $|\applytrans{\delta}{A}{u}| = 1$
	and $|\applytrans{\delta}{B}{u}| > 1$.
	
	$\ref{dist:item:subsets} \Rightarrow \ref{dist:item:subset_single_point}$:
	For $A \subseteq Q$ and $q \in Q \setminus A$ the claim follows by applying Item~\ref{dist:item:subsets} to~$A$ and $B = A \cup \{q\}$.	
	
	$\ref{dist:item:subset_single_point} \Rightarrow \ref{dist:item:single_point}$: 
	For $q \in Q$ the claim follows	by applying Item~\ref{dist:item:subset_single_point} to $A = Q \setminus \{q\}$ and $q \in Q \setminus A$.
		
	$\ref{dist:item:single_point} \Rightarrow \ref{dist:item:one-two-set}$:
	Let $A,B \subseteq Q$ be two distinct subsets of size at least two.
	Without loss of generality, suppose $q \in A \setminus B$.
	By assumption, there exists $u \in \Sigma^*$
	such that $|\delta(Q \setminus \{q\}, u)| = 1$
	and $\delta(q, u) \notin \delta(Q \setminus \{q\})$.
	Since $B \subseteq Q \setminus\{q\}$
	and $A$ contains some state distinct from~$q$,    
	we have $|\delta(B, u)| = 1$
	and $|\delta(A, u)| = 2$.	
	
	$\ref{dist:item:one-two-set} \Rightarrow \ref{dist:item:pqr-r-one-lane}$:
	Let $p,q,r\in Q$ be distinct. Consider
	$Q \setminus \{r\}$ and $Q$.
	By assumption, there exists a word $u \in \Sigma^*$ mapping one to a singleton
	and the other to a $2$-set,
	where we must have $|\delta(Q \setminus \{r\}, u)| = 1$
	and $|\delta(Q, u)| = 2$, which gives $\delta(r, u) \notin \delta(Q \setminus \{r\}, u)$.
	Then as $p,q \in Q \setminus\{r\}$ we have $\delta(p, u) = \delta(q, u)$
	and, as for $s \in Q \setminus \{r\}$ we have $\delta(s, u) = \delta(p, u)$,
	we have $\delta^{-1}(\{\delta(r,u)\}, u) = \{ s \in Q \mid \delta(s, u) = \delta(r, u)\} = \{ r \}$, i.e.,
	$r$ is the only state mapped to $\delta(r, u)$.	
		
	$\ref{dist:item:pqr-r-one-lane} \Rightarrow \ref{dist:item:compress-two-sets}$: Pick distinct states $p,q \in A$ and $r \in B \setminus A$.
	By assumption there exists $u \in \Sigma^*$
	such that $\delta(p, u) = \delta(q, u)$
	and $\delta(r, u) \ne \delta(s, u)$ for all $s \in Q \setminus \{r\}$.
	So $\delta(r, u) \notin \delta(A, u)$
	and thus $|\delta(A, u)| < |\delta(B, u)|$.
	Furthermore, since $\{p,q\} \subseteq B$
	we have $|\delta(B, u)| < |B|$.
	
	$\ref{dist:item:compress-two-sets} \Rightarrow \ref{dist:item:compress}$: 
	Consider~$A$ and $B = A \cup \{p\}$.
	By assumption, we have $|\delta(A, u)| < |\delta(B, u)| < |B|$
	for some $u \in \Sigma^*$.
	Since $|\delta(A, u)| < |\delta(B, u)|$,
	we find $\delta(p, u) \notin \delta(A, u)$.	
	We have $|B| = |A| + 1$ and $\delta(B, u) = \delta(A, u) \cup \delta(\{p\}, u)$
	so that $|\delta(B, u)| = |\delta(A, u)| + 1$.
	Thus,  $|\delta(B, u)| < |B|$
	yields $|\delta(A, u)| < |A|$.
			
	$\ref{dist:item:compress} \Rightarrow \ref{dist:item:dcps-nondisjoint}$:
	Let $A, B \subseteq Q$ be distinct with $|A|, |B| \ge 2$ and $A \cap B \ne \emptyset$.
	Without loss of generality, suppose $p \in B \setminus A$ for some $p \in Q$.
	By assumption we have $|\delta(A, u)| < |A|$
	and $\delta(p, u) \notin \delta(A, u)$ for some $u \in \Sigma^*$.
	Then we get $|\applytrans{\delta}{A}{u}| + |\applytrans{\delta}{B}{u}| < |A| + |B|$
	and, as $\delta(p, u) \in \delta(B, u)$, we get
	$\delta(A, u) \ne \delta(B, u)$.
	Furthermore, $\delta(A \cap B, u) \cup \{\delta(p,u)\} \subseteq \delta(B, u)$ with $A \cap B \ne \emptyset$ and $\delta(p, u) \notin \delta(A \cap B, u)$ gives $|\delta(B, u)| \ge 2$ so that $|\applytrans{\delta}{A}{u}| + |\applytrans{\delta}{B}{u}| \ge 3$.
	
	$\ref{dist:item:dcps-nondisjoint} \Rightarrow \ref{dist:item:dcps}$: Let $A, B \subseteq Q$ be distinct with $|A|, |B| \ge 2$. Without loss of generality, let $p \in B \setminus A$ and $q \in B \setminus\{p\}$. By the assumption applied to the distinct non-empty and non-singleton sets $A \cup \{q\}$ and $A \cup \{p,q\}$ there exists $u \in \Sigma^*$ with $3 \le |\delta(A \cup \{q\}, u)| + |\delta(A \cup \{p, q\}, u)| < 2|A|+3$ (where we even have 
	$|\delta(A \cup \{q\}, u)| + |\delta(A \cup \{p, q\}, u)| < 2|A| + 1$ for $q \in A$) and $\delta(A \cup \{q\}, u) \ne \delta(A \cup \{p, q\}, u)$. Since $\delta(A \cup \{p, q\}, u) = \delta(A \cup \{q\}, u) \cup \delta(\{p\}, u)$ we find $\delta(p, u) \notin \delta(A \cup \{q\}, u)$.
	Hence, as $\delta(p, u) \in \delta(B, u)$, we have $\delta(A, u) \ne \delta(B, u)$ 
	and $|\delta(B, u)| \ge 2$ as $\{\delta(p,u), \delta(q, u)\} \subseteq \delta(B, u)$
	where $\delta(p, u) \ne \delta(q, u)$,	
	which gives $3 \le |\delta(A, u)| + |\delta(B, u)|$. Next, for $\delta(\{p,q\}, u) \cap \delta(A, u) = \emptyset$ we have 
	\begin{align*}
		2|\delta(A, u)| + 3 & = |\delta(A, u)| + |\delta(\{q\}, u)| + |\delta(A, u)| + |\delta(\{p,q\}, u)| \\
		& = |\delta(A \cup \{q\}, u)| + |\delta(A \cup \{p, q\}, u)| \\
		& < 2|A| + 3.
	\end{align*}
	Now, $2|\delta(A, u)| + 3 < 2|A| + 3$ gives $|\delta(A, u)| < |A|$. So, we find $|\delta(A, u)| + |\delta(B, u)| < |A| + |B|$. Otherwise, suppose $\delta(\{p,q\}, u) \cap \delta(A, u) \ne \emptyset$, which gives $\delta(A, u) \cap \delta(B, u) \ne \emptyset$. In this case, either we have $|\delta(A, u)| + |\delta(B, u)| < |A| + |B|$ and the claim holds with~$u$, or $|\delta(A, u)| = |A| \ge 2$ and $|\delta(B, u)| = |B| \ge 2$, and by the assumption applied to $\delta(A, u)$ and $\delta(B, u)$ there exists $v \in \Sigma^*$ with $3 \le |\delta(A, uv)| + |\delta(B, uv)| < |\delta(A, u)| + |\delta(B, u)|$ and $\delta(A, uv) \ne \delta(B, uv)$. Then the claim holds with~$uv$ as $|\delta(A, u)| + |\delta(B, u)| \le |A| + |B|$.
	
	$\ref{dist:item:dcps} \Rightarrow \ref{dist:item:all_pairs}$:
	Let $A,B \subseteq Q$
	be distinct with $|A|, |B| \ge 2$.
	We do induction on $|A| + |B|$.
	First, suppose $|A| + |B| = 4$, which is the base case.
	By assumption, there exists $u \in \Sigma^*$
	such that $3 \le |\applytrans{\delta}{A}{u}| + |\applytrans{\delta}{B}{u}|  < 4$, which
	implies that one of both sets must be a singleton and the other a $2$-set.
	Hence, $u \in \Sigma^*$ distinguishes both subsets.
	Next, suppose $|A| + |B| > 4$. Assume all pairs of distinct subsets $A', B' \subseteq Q$
	with $4 \le |A'| + |B'| < |A| + |B|$
	and $|A'|, |B'| \ge 2$ are distinguishable.
	By the assumption from~\cref{dist:item:dcps} there exists
	$u \in \Sigma^*$
	such that  $3 \le |\applytrans{\delta}{A}{u}| + |\applytrans{\delta}{B}{u}|  < |A| + |B|$
	and $\applytrans{\delta}{A}{u} \ne \applytrans{\delta}{B}{u}$.
	If one of both subsets is a singleton, the other subset has size at least two.
	So, $u$ distinguishes~$A$ and~$B$.
	Otherwise, as $\applytrans{\delta}{A}{u} \ne \applytrans{\delta}{B}{u}$, by induction hypothesis, 
	there exists $v \in \Sigma^*$
	distinguishing $\applytrans{\delta}{A}{u}$ from $\applytrans{\delta}{B}{u}$.
	Then~$uv$ distinguishes~$A$ from~$B$.
\end{proof}

\Crefitem{sync_max:thm:char_distinguishable}{dist:item:dcps} is useful for an inductive reasoning
because of the fact that it only demands to compress the subsets (while retaining their distinctiveness)
instead of instantly mapping precisely one to a singleton subset.

\begin{remark}
 By \cref{sync_max:ex:distinguishable-trap} 	
 distinguishability of all pairs of distinct $(n-1)$-sets 
 is not sufficient for complete distinguishability.
 With \cref{sync_max:prop:ksets},
 this also implies that there exists no~$k$
 such that distinguishability of all pairs of distinct $k$-sets
 implies complete distinguishability.
\end{remark}

\begin{remark}
  Having for each $q \in Q$ a word $u_q \in \Sigma^*$
  that distinguishes $Q$ form $Q \setminus \{q\}$
  is necessary, and, by \crefitem{sync_max:thm:char_distinguishable}{dist:item:single_point}
  also sufficient, for complete distinguishability.
  Let $p,q,r \in Q$ be distinct.
  Assume $\delta(r, u_q) = \delta(r, u_p)$.
  Then we have $\delta(p, u_q) = \delta(r, u_q) = \delta(r, u_p) = \delta(q, u_p)$.
  Since $\delta(q, u_q) \ne \delta(p, u_q)$,
  we find  $\delta(q, u_q) \ne \delta(q, u_p)$.
  Thus we conclude that each word $u_q$ induces a different
  transformation on the states.
  Hence, the transformation monoid of completely distinguishable
  automata with at least three states contains at least~$|Q|$ different elements.
\end{remark}

\begin{corollary}\label{sync_max:cor:sufficient-compl-dist}
	Let $\semiautomat{A}{Q}{\Sigma}{\delta}$.
	If for all $r \in Q$
	and $p, q \in Q \setminus \{q\}$
	there exists $u \in \Sigma^*$
	such that $\delta(p, u) = \delta(q, u)$, $\delta(r, u) = r$,
	and $\delta(s, u) \ne r$ for all $s \in Q \setminus \{r\}$, then $\mathcal A$
	is completely distinguishable.
\end{corollary}
\begin{proof}
	This follows directly from \crefitem{sync_max:thm:char_distinguishable}{dist:item:pqr-r-one-lane}
	since it equals this statement in the special case that the word fixes the state~$r$.
\end{proof}

\begin{corollary}\label{cor:compl-dist-transitive}
	Let $\semiautomat{A}{Q}{\Sigma}{\delta}$
	be an automaton such that the group of units of~$\TM{\mathcal A}$
	is transitive.
	Then $\mathcal A$ is completely distinguishable
	if and only if $Q$ is distinguishable from $Q \setminus \{q\}$
	for some $q \in Q$.		
\end{corollary}
\begin{proof}
	First, suppose $Q$ and $Q \setminus \{q\}$ are distinguishable for $q \in Q$.
	Then $|\delta(Q \setminus \{q\}, u)| = 1$
	with $\delta(q, u) \notin \delta(Q \setminus\{q\}, u)$
	for some $u \in \Sigma^*$.
	Let $p \in Q$.
	Since~$\TM{\mathcal A}$ contains a transitive permutation group, 
	there exists $v \in \Sigma^*$ inducing a permutation on $Q$
	such that $\delta(p, v) = q$.
	Then we have $\applytrans{\delta}{Q \setminus \{p\}}{vu} = \applytrans{\delta}{Q\setminus\{q\}}{u}$
	and $\applytrans{\delta}{p}{vu} = \applytrans{\delta}{q}{u}$.
	So $\applytrans{\delta}{p}{vu} \notin \applytrans{\delta}{Q \setminus\{p\}}{vu}$
	and $|\applytrans{\delta}{Q\setminus \{p\}}{vu}| = 1$.
	By~\crefitem{sync_max:thm:char_distinguishable}{dist:item:single_point},
	it follows that $\mathcal A$ is completely distinguishable.	
	The converse implication is clear.
\end{proof}

\begin{example}  \label{sync_max:ex:compl-dist-non-cr}
	Consider the automaton
	in \cref{sync_max:fig:compl_dist_not_compl_reach:1}.
	The letter $a$ distinguishes $\{0,1,2,3\}$
	from the whole state set. 
	Since the letter $b$ generates a transitive permutation group,
	the automaton is completely distinguishable by \cref{cor:compl-dist-transitive}.	
	The automaton is not completely reachable
	since it does not have a letter of defect one.
\end{example}

\begin{example} \label{sync_max:ex:compl-dist-non-cr2}
	\Cref{sync_max:fig:compl_dist_not_compl_reach:2}
	depicts an automaton that is easily seen
	to be completely distinguishable
	and that is not completely reachable
	since it is not strongly connected.
\end{example}

\begin{example} \label{syncmax:ex:compl-dist-not-SCC}
	The word $ab$ distinguishes $\{1,2,3\}$ and $\{4\}$.
	The word $ecab$ distinguishes $\{2,3,4\}$ and $\{1\}$.
	The word $bc$ distinguishes $\{1,2,4\}$ and $\{3\}$.
	The word $ac$ distinguishes $\{1,3,4\}$ and $\{2\}$.
	So, by \crefitem{sync_max:thm:char_distinguishable}{dist:item:single_point},
	the automaton from \cref{sync_max:fig:compl_dist_example} is completely distinguishable,
	but it is not completely reachable since it is not strongly connected.
	Note that by adding a simple idempotent mapping $1$ to $2$, we obtain
	a strongly connected SI-automaton that is completely distinguishable,
	but whose automaton graph is not strongly $\frac{3}{2}$-connected
	since all paths from $2$ to $4$ have to pass through~$1$.
\end{example}

\begin{figure}[ht]
	\centering 
	\begin{subfigure}[t]{0.37\textwidth}
		\centering
		\begin{tikzpicture}[shorten >=1pt,->,node distance=4cm]
				\foreach \name/\angle/\text in {P-n-2/234/3, P-n-1/162/4, 
					P-0/90/0, P-1/18/1, P-2/-54/2}
				\node[state] (\name) at (\angle:7em) {$\text$};
				
				\path[->] (P-n-2) edge [left] node {$b$} (P-n-1);
				\path[->] (P-n-1) edge [above,pos=0.35] node {$b$} (P-0);
				\path[->] (P-0) edge [above] node {$b$} (P-1);
				\path[->] (P-1) edge [right] node {$b$} (P-2);
				\path[->] (P-2) edge [below] node {$a,b$} (P-n-2);
				\path[->] (P-0)   edge [right] node {$a$} (P-n-2)
		                  (P-1)   edge [above] node {$a$} (P-n-2);					
			\end{tikzpicture}
		\caption{The Strongly connected automaton from Example~\ref{sync_max:ex:compl-dist-non-cr}.}
		\label{sync_max:fig:compl_dist_not_compl_reach:1}
	\end{subfigure}\hfill
	\begin{subfigure}[t]{0.29\textwidth}
		\centering		
		\begin{tikzpicture}
			\node[state] (1) {$1$};
			\node[state, above = of 1] (2) {$2$};
			\node[state, left  = of 2] (3) {$3$};    		
			
			\path[->] (3) edge [left]  node {$a$} (1)
			          (2) edge [right] node {$b$} (1)
			          (2) edge [above] node {$c$} (3);
		\end{tikzpicture}		
		\caption{The automaton from Example~\ref{sync_max:ex:compl-dist-non-cr2}.} 
		\label{sync_max:fig:compl_dist_not_compl_reach:2}
	\end{subfigure}\hfill
    \begin{subfigure}[t]{0.29\textwidth}
    	\centering
    		\begin{tikzpicture}[node distance = 3.5em and 1.5em]
    		\node[state] (1) {$1$};
    		\node[state, above right = of 1] (2) {$2$};
    		\node[state, above left  = of 1] (3) {$3$};
    		\node[state, below = of 1]       (4) {$4$};				
    		
    		\path[->] (3) edge [left]  node {$a$} (1)
    		          (2) edge [right] node {$b$} (1)
    		          (1) edge [bend left=10, right] node {$c$} (4)
    		          (4) edge [bend left=10, left]  node {$d$} (1)
    		          (4) edge [bend left, left]  node {$e$} (3);
    	\end{tikzpicture}
        \caption{The automaton from \cref{syncmax:ex:compl-dist-not-SCC}}     	
        \label{sync_max:fig:compl_dist_example}    
    \end{subfigure}	\caption[Completely distinguishable, but not completely reachable]{Completely 
    distinguishable automata that are not completely reachable. Self-loops omitted.}
	\label{sync_max:fig:compl_dist_not_compl_reach}
\end{figure}

It is easily seen that quotient automata of completely reachable automata
are completely reachable. However, complete distinguishability is not preserved
under taking quotient automata as shown in \cref{sync_max:fig:compl_dist_not_compl_reach} next.

% Und sync-maximal erhalten unter Quotientautomaten?
\begin{example} \label{sync_max:fig:compl_dist_not_compl_reach}
	Automata  $\semiautomat{A}{Q}{\Sigma}{\delta}$ as in \cref{sync_max:ex:compl-dist-non-cr}
	also show that complete distinguishability
	is not preserved when taking quotient automata.
	For consider the obvious generalization	where the letter~$b$ cyclically permutes $n$ states
	and suppose we have some non-trivial $\TS{\mathcal A}$-congruence $\sim$.
	%Let $p \in Q$ be the state fixed by~$a$.
	Let $s \in Q$ be the state fixed by~$a$.
	Since the letter $b$ cyclically permutes the whole state set,
	all congruence classes for~$\sim$ have the same size
	and we find $t \in Q$ with $s \ne t$
	such that $s \sim t$.
		
	First, we show that for every $u \in \Sigma^*$ not permuting the set of states
	we have $\delta(p, u) \sim \delta(q, u)$
	for all $p,q \in Q$, i.e., the letter~$a$ induces a constant in $\mathcal A / {\sim}$).
	Let $p,q \in Q$. We can assume $p \not\sim q$.
	For $u \in \Sigma^*$ not permuting~$Q$
	we can write $u = b^iav$ for some $v \in \Sigma^*$ and $i \ge 0$.
	If $\delta(p, b^i) \ne s$ 
	and $\delta(q, b^i) \ne s$,
	then $\delta(p, b^ia) = \delta(q, b^ia)$
	and so $\delta(p, u) = \delta(q, u)$.
	Otherwise, without loss of generality, assume $\delta(p, b^i) = s$.
	Then $\delta(q, b^i) \ne s$.
	Now, we have $\delta(p, b^ia) = s$
	and $\delta(q, b^ia) = \delta(t, a)$.
	Furthermore, $s \sim t$
	implies $\delta(p, b^ia) = s = \delta(s, a) \sim \delta(t, a) = \delta(q, b^ia)$
	and thus $\delta(p, u) \sim \delta(q, u)$.
		
	Next, assume we want to distinguish $Q / {\sim}$
	and $Q / {\sim} \setminus \{ [p]_{\sim} \}$
	in $\mathcal A / {\sim}$ for some $p \in Q$.
	Then a distinguishing word $u$
	maps all $q \in Q$ with $q \not\sim p$
	into an equivalence class
	and for all $r \in Q$ with $r \sim p$
	we have $\delta(r, u) \not\sim \delta(q, u)$ for $q \in Q$ with $q \not\sim p$.
	However, since a distinguishing word cannot permute
	the whole state set, as shown above this is not possible.	
\end{example}

Next, we show that a binary completely distinguishable automaton with a letter of rank two, such as the one from \cref{sync_max:fig:compl_dist_not_compl_reach:1}, must contain a letter cyclically permuting the states.

\begin{proposition} \label{syncmax:prop:binary-rank-two-letter}
	Let $\semiautomat{A}{Q}{\{a,b\}}{\delta}$ be a binary automaton with at least three states
	where~$a$ has rank two.
	If $\mathcal A$ is completely distinguishable,
	then~$b$ cyclically permutes the states.
\end{proposition}
\begin{proof}
	Assume $\delta(p_1, b) = \delta(p_2, b)$
	for two distinct states $p_1, p_2 \in Q$.
	Then a word distinguishing $\{p_1\}$ from $Q \setminus\{p_1\}$
	must start with~$a$.
	This implies $\delta(p_1, a) \notin \delta(Q \setminus \{p_1\}, a)$.
	Since~$a$ has rank two,
	we get $|\delta(Q \setminus \{p_1\}, a)| = 1$
	and so~$a$ distinguishes both subsets.
	By a similar reasoning with $\{p_2\}$ and $Q \setminus \{p_2\}$,
	we find $|\delta(Q \setminus \{p_2\}, a)| = 1$.
	Since~$|Q| \ge 3$,
	this yields that $Q \setminus \{p_1\}$
	and $Q \setminus \{p_2\}$ are mapped to the same
	state and so $\delta(p_1, a) = \delta(p_2, a)$,
	giving a contradiction
	to $\delta(p_1, a) \notin \delta(Q \setminus \{p_1\}, a)$.
	Thus, no two distinct states are mapped
	to the same state by~$b$,
	and by finiteness we find that~$b$
	permutes the states.
	
	Next, by assumption we have $Q = A \cup B$
	where $A = \delta^{-1}(p, a)$
	and $B = \delta(q, a)$ for some $p,q \in Q$.
	Assume there exist two distinct $|A|, |B| \ge 2$.
	Each word distinguishing two subsets
	must have the form $b^ia$ for some $i \ge 0$,
	since~$a$ has rank two.
	We have $\delta^{-1}(p, b^ia) = \delta^{-1}(A, b^i)$
	and $\delta^{-1}(q, b^ia) = \delta^{-1}(B, b^i)$
	and the two subsets have size at least two
	and partition~$Q$ since~$b$ permutes the states.
	However, a word distinguishes a state~$s$
	from $Q \setminus \{s\}$
	if and only if $\{s\}$
	and $Q \setminus \{s\}$
	are the $\delta^{-1}$-images of two states.
	Thus, one of both subsets $A$ or $B$
	must be a singleton.
	Without loss of generality, let~$A = \{p_1\}$ and consider some state $p_2 \in Q$.
	A word $b^ia$ distinguishes $p_2$ and $Q \setminus \{p_2\}$
	if and only if $p_2$ is mapped to~$p_1$
	by~$b^i$. Thus we find that~$b$ cyclically permutes the states.	
\end{proof}

Lastly, we also note a readily seen equivalence when the whole state set is distinguishable
from a set of the form $Q \setminus\{q\}$, since these sets appear in \crefitem{sync_max:thm:char_distinguishable}{dist:item:single_point}

% in beweis oben nutzen?
\begin{proposition} \label{sync_max:prop:dist-n-1-set}
	Let $\semiautomat{A}{Q}{\Sigma}{\delta}$. Then the following are equivalent.
	\begin{enumerate}
		\item $Q \setminus \{q\}$ is distinguishable from~$Q$.
		
		\item There exist $u_1, \ldots, u_m \in \Sigma$
		such that for  $A_i = \delta(Q \setminus \{q\}, u_1 \cdots u_i)$
		and $q_i = \delta(q, u_1 \cdots u_i)$ where $i \in \{1,\ldots,m\}$
		we have  $|A_m| = 1$
		and $q_i \notin A_i$ for all $i \in \{1,\ldots,m\}$.
	\end{enumerate}
\end{proposition}

\begin{example}
	Let $\semiautomat{A}{\{1,2,3,4\}}{\Sigma}{\delta}$
	with 
	\[
	\Sigma = \{ a \mid \exists y_1, y_2 \in Q \colon y_1 \ne y_2 \land
	|\delta^{-1}(y_1, a)| = |\delta^{-1}(y_1, a)| = 2 \}
	\]
	Then disjoint pairs of $2$-sets are indistinguishable
	and non-disjoint pairs of $2$-sets are distinguishable.
	Furthermore, two sets of size at least three are indistinguishable.
	A set of size two is distinguishable from any set of size at least three.	
\end{example}